\section*{Примеры}
\addcontentsline{toc}{section}{Примеры}

\subsection*{Пример 1. Как из линейной оболочки (линейного подпространства, натянутого на систему векторов) получить систему уравнений}

Пусть
\[
L=
\left\langle
\begin{pmatrix}
1\\0\\1\\1
\end{pmatrix},
\begin{pmatrix}
0\\1\\1\\0
\end{pmatrix}
\right\rangle
\subset \mathbb{R}^4.
\]
Требуется получить систему линейных уравнений, задающую это подпространство.

\textbf{Идея решения.}

Составим матрицу из образующих векторов, записав их строками:
\[
E=
\begin{pmatrix}
1&0&1&1\\
0&1&1&0
\end{pmatrix}.
\]

Далее рассматриваем однородную систему
\[
Ey=0.
\]
После приведения к ступенчатому виду выделяются главные и свободные переменные, затем строится фундаментальная система решений. Именно она и позволяет записать искомое подпространство в виде системы уравнений.

\medskip
\noindent
\textbf{Комментарий к алгоритму.}
Если подпространство задано как линейная оболочка векторов, то удобно перейти к однородной системе и найти её ФСР. Из неё затем восстанавливается система линейных уравнений, задающая это подпространство.

\subsection*{Пример 2. Как находить базис и размерность линейной оболочки}

Пусть в $\mathbb R^4$ задана система векторов
\[
a_1=
\begin{pmatrix}
1\\0\\1\\0
\end{pmatrix},
\quad
a_2=
\begin{pmatrix}
0\\1\\1\\1
\end{pmatrix},
\quad
a_3=
\begin{pmatrix}
1\\1\\2\\1
\end{pmatrix},
\quad
a_4=
\begin{pmatrix}
2\\1\\3\\1
\end{pmatrix}.
\]
Требуется найти базис и размерность линейной оболочки этих векторов.

\textbf{Идея решения.}

Записываем векторы строками матрицы и приводим её к ступенчатому виду методом Гаусса. Число ненулевых строк после приведения совпадает с размерностью линейной оболочки. В качестве базиса удобно брать те из исходных векторов, которые соответствуют ведущим строкам.

\medskip
\noindent
\textbf{Комментарий к алгоритму.}
Для задач на базис и размерность линейной оболочки самый надёжный путь — ранговый: составить матрицу из данных векторов и сделать элементарные преобразования строк.

\subsection*{Пример 3. Как находить сумму и пересечение подпространств}

Пусть
\[
U=
\left\langle
\begin{pmatrix}
1\\0\\1\\0
\end{pmatrix},
\begin{pmatrix}
0\\1\\1\\0
\end{pmatrix}
\right\rangle,
\qquad
V=
\left\langle
\begin{pmatrix}
1\\1\\2\\1
\end{pmatrix},
\begin{pmatrix}
0\\0\\0\\1
\end{pmatrix}
\right\rangle
\subset \mathbb R^4.
\]
Требуется найти $U\cap V$, $U+V$ и их размерности.

\textbf{Идея решения.}

Сначала найдём размерности самих подпространств по рангу матрицы, составленной из их образующих. Затем объединим все образующие в одну матрицу и найдём размерность суммы. После этого используем формулу Грассмана:
\[
\dim(U+V)=\dim U+\dim V-\dim(U\cap V).
\]
Если требуется базис суммы, его удобно взять из линейно независимых векторов объединённой системы образующих.

\medskip
\noindent
\textbf{Комментарий к алгоритму.}
Для таких задач удобно работать через ранги и формулу Грассмана: сначала находят размерности подпространств и их суммы, а затем из формулы получают размерность пересечения.

\subsection*{Пример 4. Прямая сумма и проекции}

Рассмотрим пространство $\mathbb R^3$ и подпространства
\[
L_1=\left\{(x_1,x_2,x_3)\in\mathbb R^3 \mid x_1+x_2+x_3=0\right\},
\qquad
L_2=\left\langle
\begin{pmatrix}
1\\1\\1
\end{pmatrix}
\right\rangle.
\]
Требуется доказать, что
\[
\mathbb R^3=L_1\oplus L_2,
\]
и найти проекции вектора
\[
x=
\begin{pmatrix}
2\\-1\\3
\end{pmatrix}
\]
на $L_1$ параллельно $L_2$ и на $L_2$ параллельно $L_1$.

\textbf{Идея решения.}

Сначала проверяем, что пересечение $L_1\cap L_2$ тривиально. Затем сравниваем суммы размерностей: если
\[
\dim L_1+\dim L_2=\dim \mathbb R^3
\quad\text{и}\quad
L_1\cap L_2=\{0\},
\]
то получаем прямую сумму.

После этого вектор $x$ раскладывается в виде
\[
x=u+v,\qquad u\in L_1,\ v\in L_2.
\]
Неизвестный коэффициент при векторе из $L_2$ находится из условия принадлежности $u$ подпространству $L_1$. Полученное разложение и даёт обе проекции.

\medskip
\noindent
\textbf{Комментарий к алгоритму.}
Во всех задачах на прямую сумму удобно сначала доказать её существование, а уже потом искать разложение вектора по компонентам.

\subsection*{Пример 5. Разложение матрицы на симметрическую и кососимметрическую части}

Пусть
\[
A=
\begin{pmatrix}
1&2&0\\
0&3&4\\
5&1&2
\end{pmatrix}.
\]
Требуется разложить матрицу $A$ в сумму симметрической и кососимметрической матриц.

\textbf{Идея решения.}

Для любой квадратной матрицы используется стандартная формула
\[
A=\frac{A+A^T}{2}+\frac{A-A^T}{2}.
\]
Первая матрица симметрична, вторая кососимметрична. Поэтому задача сводится к вычислению транспонированной матрицы и двух простых выражений.

\medskip
\noindent
\textbf{Комментарий к алгоритму.}
Это универсальный способ для разложения любой квадратной матрицы на симметрическую и кососимметрическую части.

\subsection*{Пример 6. Как из системы уравнений получить линейную оболочку}

Пусть подпространство $L \subset \mathbb R^4$ задано системой уравнений
\[
\begin{cases}
x_1+x_2-x_3=0,\\
x_1-x_4=0.
\end{cases}
\]
Требуется найти базис и размерность подпространства $L$ и записать его в виде линейной оболочки.

\textbf{Идея решения.}

Рассматриваем это подпространство как множество решений однородной системы. Затем выражаем главные переменные через свободные, вводим параметры и строим фундаментальную систему решений. Векторы ФСР образуют базис пространства решений, а само пространство записывается как их линейная оболочка.

\medskip
\noindent
\textbf{Комментарий к алгоритму.}
Если подпространство задано системой $Ax=0$, то нужно найти ФСР. Это сразу даёт и базис, и размерность.

\medskip
\noindent
\textbf{Итог.}
Во всех примерах мы используем одни и те же базовые приёмы: метод Гаусса, параметризацию, формулу Грассмана, критерий прямой суммы и стандартное разложение матрицы на симметрическую и кососимметрическую части.
